\documentclass[aspectratio=169]{beamer}
\input{header}
\coursecode{JKSSB}

\title{Access Database Basics}
\subtitle{Lesson 28 --- Databases, Tables and Records}
\author{JKSSB Computer Awareness}
\date{\today}

\begin{document}
\frame{\titlepage}

\begin{frame}{What Is a Database?}
  \begin{itemize}
    \item A \key{database} is an organised collection of related data.
    \item The data is stored so that it can be \key{searched, updated, and
      retrieved} easily.
    \item Example: a school's list of students, their classes, and their
      marks kept together as related information.
    \item In a computer, a database is managed by special software, not kept
      as loose files.
  \end{itemize}
  \begin{block}{The one idea}
    A database is the \key{collection of data}; the software that manages it
    is a separate thing.
  \end{block}
\end{frame}

\begin{frame}{Why Use a Database Instead of a Spreadsheet?}
  \begin{itemize}
    \item A spreadsheet is built for \muted{calculation} on a simple grid.
    \item A database is built for \key{storing and linking large amounts of
      related data}.
    \item A database can hold many tables, relate them, and answer queries
      across them.
    \item For a big, growing record system, a \key{database} keeps data
      consistent and avoids duplication.
  \end{itemize}
  \begin{alertblock}{Keep it clear}
    Excel is a spreadsheet; Access is a database program. The two are not
    the same tool.
  \end{alertblock}
\end{frame}

\begin{frame}{Database vs DBMS (Near-Neighbour)}
  \begin{center}
    \begin{tabular}{p{0.20\textwidth}p{0.33\textwidth}p{0.37\textwidth}}
      \toprule
      \textbf{Point} & \textbf{Database} & \textbf{DBMS} \\
      \midrule
      What it is & The organised collection of data & The software that manages
        the data \\
      Nature & Data & Program \\
      Example & A school's records & MS Access, Oracle, MySQL \\
      \bottomrule
    \end{tabular}
  \end{center}
  \vspace{0.25cm}
  \muted{The database is the data; the DBMS is the tool. Do not swap them.}
\end{frame}

\begin{frame}{What a DBMS Does}
  \begin{itemize}
    \item \key{DBMS} stands for \key{Database Management System}.
    \item It \key{stores} data and controls how it is kept on disk.
    \item It lets users \key{insert, update, delete, and retrieve} records.
    \item It \key{controls access} so that only allowed users can read or
      change data.
    \item It keeps the data \key{consistent} and safe from corruption.
  \end{itemize}
  \begin{block}{Full form to memorise}
    DBMS = Database Management System. An \key{RDBMS} is a relational DBMS
    (data stored in related tables).
  \end{block}
\end{frame}

\begin{frame}{Microsoft Access in the Office Suite}
  \begin{itemize}
    \item \key{Microsoft Access} is the database program of the
      \key{Microsoft Office} suite.
    \item It is a \key{relational database management system} (RDBMS).
    \item In the suite, Word handles text, Excel handles spreadsheets,
      PowerPoint handles presentations, and \key{Access handles databases}.
    \item Access stores tables, queries, forms, and reports inside one
      database file.
  \end{itemize}
  \begin{alertblock}{Exam point}
    Access is a DBMS, not a spreadsheet. It is part of the MS Office
    family.
  \end{alertblock}
\end{frame}

\begin{frame}{The Objects in an Access Database}
  \begin{center}
    \begin{tabular}{p{0.18\textwidth}p{0.68\textwidth}}
      \toprule
      \textbf{Object} & \textbf{Job} \\
      \midrule
      Table & Stores the actual data in rows and columns \\
      Query & Asks a question and retrieves selected data \\
      Form & Provides a screen to enter and view data \\
      Report & Presents data in a printable layout \\
      \bottomrule
    \end{tabular}
  \end{center}
  \vspace{0.25cm}
  \muted{Data lives in tables; the other objects work on that data.}
\end{frame}

\begin{frame}{The Table: Rows and Columns}
  \begin{itemize}
    \item A \key{table} is the basic place where data is stored.
    \item Data is arranged in \key{rows and columns}.
    \item Each \key{column} is one \key{field} --- one kind of information.
    \item Each \key{row} is one \key{record} --- one complete entry.
    \item A table therefore holds \key{many records}, each with the same set
      of fields.
  \end{itemize}
  \begin{block}{The building block}
    Everything in an Access database starts with the \key{table}.
  \end{block}
\end{frame}

\begin{frame}{Field vs Record vs Table (Near-Neighbour)}
  \begin{center}
    \begin{tabular}{p{0.16\textwidth}p{0.34\textwidth}p{0.30\textwidth}}
      \toprule
      \textbf{Term} & \textbf{Meaning} & \textbf{Position in a table} \\
      \midrule
      Field & One kind of information & A column \\
      Record & One complete entry & A row \\
      Table & A collection of related records & The whole grid \\
      \bottomrule
    \end{tabular}
  \end{center}
  \vspace{0.25cm}
  \begin{alertblock}{Trap alert}
    Do not mix these up: a \key{field} is a column, a \key{record} is a
    \key{row}, and a \key{table} is the whole collection.
  \end{alertblock}
\end{frame}

\begin{frame}{Data Types in Access}
  \begin{itemize}
    \item Every field has a \key{data type} --- the kind of value it can
      hold.
    \item \key{Short Text} --- names, codes, and other text.
    \item \key{Number} --- values used in calculations.
    \item \key{Date/Time} --- dates and times.
    \item \key{Yes/No} --- only True or False (a Boolean value).
    \item \key{AutoNumber} --- a unique number added automatically for each
      new record.
  \end{itemize}
  \begin{block}{Why it matters}
    A field can hold only values that match its data type, which keeps the
    data clean.
  \end{block}
\end{frame}

\begin{frame}{Field Properties}
  \begin{itemize}
    \item A \key{field property} is a setting that describes how a field
      behaves.
    \item \key{Field Size} limits how much a field can hold.
    \item \key{Default Value} fills in a value automatically.
    \item \key{Required} --- if set to Yes, the field \key{cannot be left
      blank}.
    \item \key{Input Mask} fixes the pattern of entry, such as a phone
      number.
  \end{itemize}
\end{frame}

\begin{frame}{Input Validation}
  \begin{itemize}
    \item \key{Validation} means checking that entered data is sensible and
      allowed.
    \item A \key{Validation Rule} states the condition the value must meet.
    \item Example rule: \key{Between 0 And 100} --- only values in that range
      are accepted.
    \item \key{Validation Text} is the message shown when the rule is
      violated.
  \end{itemize}
  \begin{exampleblock}{Worked example}
    Set a Validation Rule of \texttt{Between 0 And 100} on a Marks field.
    If someone types 150, Access rejects it and shows the Validation Text.
  \end{exampleblock}
\end{frame}

\begin{frame}{The Primary Key}
  \begin{itemize}
    \item A \key{primary key} is a field that \key{uniquely identifies} each
      record.
    \item Its value is \key{different for every record} and is never left
      blank.
    \item Example: a \key{Roll Number} identifies each student in a class.
    \item An \key{AutoNumber} field is often used as a ready-made primary
      key.
  \end{itemize}
  \begin{alertblock}{Trap alert}
    A primary key \key{cannot} contain duplicate values. If two records share
    the same key value, Access will not allow it.
  \end{alertblock}
\end{frame}

\begin{frame}{Indexes}
  \begin{itemize}
    \item An \key{index} is a special structure that makes searching and
      sorting a table \key{faster}.
    \item It works like the index of a book, pointing straight to the right
      place.
    \item A primary key is \key{automatically indexed}.
    \item An index takes extra space and can slightly slow down adding data,
      so it is used only where needed.
  \end{itemize}
  \begin{block}{Remember}
    Index = speed of \key{search}; it does not change the data itself.
  \end{block}
\end{frame}

\begin{frame}{Worked Example: A Students Table}
  \begin{center}
    \begin{tabular}{p{0.18\textwidth}p{0.18\textwidth}p{0.14\textwidth}p{0.16\textwidth}}
      \toprule
      \textbf{RollNo} & \textbf{Name} & \textbf{Marks} & \textbf{Passed} \\
      \midrule
      1 & Ayaan & 78 & Yes \\
      2 & Bilal & 45 & No \\
      3 & Chitra & 91 & Yes \\
      \bottomrule
    \end{tabular}
  \end{center}
  \vspace{0.2cm}
  \begin{itemize}
    \item \key{RollNo} is the primary key (unique for each student).
    \item \key{Name} is Short Text; \key{Marks} is Number; \key{Passed} is
      Yes/No.
    \item Each row is one \key{record}; each column is one \key{field}.
  \end{itemize}
\end{frame}

\begin{frame}{Trap Alert: Common Misreadings}
  \begin{itemize}
    \item A \key{database} is the data; a \key{DBMS} is the software.
      (TRAP-style)
    \item A \key{field} is a column; a \key{record} is a row.
    \item \key{Validation} checks the value; an \key{index} speeds up search.
    \item A \key{primary key} is unique; it cannot contain duplicates.
    \item Access is part of \key{MS Office}, not a spreadsheet.
  \end{itemize}
\end{frame}

\begin{frame}{Lesson 28: Recall}
  \begin{itemize}
    \item A database is an organised collection of related data; a DBMS is
      the software that manages it.
    \item Microsoft Access is the RDBMS of the MS Office suite.
    \item The table stores data in \key{fields} (columns) and \key{records}
      (rows).
    \item Data types include Short Text, Number, Date/Time, Yes/No, and
      AutoNumber.
    \item Validation Rules and Validation Text check values; Required stops
      blank entries.
    \item A primary key uniquely identifies a record; an index speeds up
      searching.
  \end{itemize}
\end{frame}
\end{document}
